\section{Extended Related Work}
\label{app:extended_related_work}

Our work extends LLM abstention research to tool-using agents. We organize related work into four areas: (1) abstention, calibration, and over-refusal in LLMs, which provides the conceptual foundation; (2) clarification and ambiguity handling, which addresses the dialogue side of restraint; (3) agent capability and tool-use benchmarks, which characterise what agents \emph{can} do; and (4) agent safety, deception, and risk benchmarks, which collectively reveal the evaluation gap our work targets.

\begin{table}[htbp]
\centering
\caption{Positioning of \benchmark{} relative to existing benchmarks. $^\dagger$Concurrent work.}
\label{tab:gaps}
\small
\begin{tabular}{lccc}
\toprule
\textbf{Benchmark} & \textbf{Tool-Use} & \textbf{Mid-Task} & \textbf{Paired} \\
 & \textbf{Env.} & \textbf{Abstention} & \textbf{Design} \\
\midrule
AbstentionBench~\citep{kirichenko2025abstentionbench} & \ding{55} & \ding{55} & \ding{55} \\
AgentHarm~\citep{andriushchenko2025agentharm} & \ding{51} & \ding{55} & \ding{55} \\
Agent-SafetyBench~\citep{zhang2025agentsafetybench} & \ding{51} & \ding{55} & \ding{55} \\
BFCL V4~\citep{patil2025bfcl} & \ding{51} & \ding{55} & \ding{55} \\
ToolBeHonest~\citep{zhang2024toolbehonest} & \ding{51} & Partial & \ding{55} \\
ToolSandbox~\citep{lu2025toolsandboxstatefulconversationalinteractive} & \ding{51} & Partial & \ding{55} \\
ToolEmu~\citep{ruan2024identifyingriskslmagents} & \ding{51} & \ding{51} & \ding{55} \\
$\tau$-bench~\citep{yao2024taubench} & \ding{51} & \ding{55} & \ding{55} \\
ST-WebAgentBench~\citep{levy2026stwebagentbench} & \ding{51} & \ding{55} & \ding{55} \\
ImpossibleBench~\citep{zhong2025impossiblebenchmeasuringllmspropensity} & \ding{51} & \ding{51} & \ding{55} \\
FixedBench~\citep{gloaguen2026codingagents} & \ding{51} & \ding{51} & \ding{55} \\
OR-Bench~\citep{cui2025orbench} & \ding{55} & \ding{55} & \ding{55} \\
Agentic Abstention~\citep{luo2026agentic}$^\dagger$ & \ding{51} & \ding{51} & \ding{55} \\
\midrule
\textbf{\benchmark{} (Ours)} & \ding{51} & \ding{51} & \ding{51} \\
\bottomrule
\end{tabular}
\end{table}

\subsection{LLM Abstention, Calibration, \& Over-Refusal}

\textbf{Abstention as a capability.}
Wen et al.~\citep{wen2025knowlimits} examine abstention from three perspectives: query-focused (is this answerable?), model-focused (can the model answer?), and human-values-focused (should it answer?). We adapt these to action contexts: task-focused, capability-focused, and consequence-focused. Madhusudhan et al.~\citep{madhusudhan2024llmsknowanswerinvestigating} introduce \textsc{Abstain-QA} along with the Answerable--Unanswerable Confusion Matrix and formalise ``Abstention Ability'' as a capability distinct from raw QA accuracy. Closely related, Tomani et al.~\citep{tomani2024uncertaintybasedabstentionllmsimproves} show that uncertainty-based abstention reduces hallucinations and improves safety, while Kapoor et al.~\citep{kapoor2025largelanguagemodelstaught} argue that LLMs must be \emph{taught} to know what they don't know via white-box calibration on labelled correctness data. Cohen et al.~\citep{cohen2024idontknowexplicit} propose an explicit \texttt{[IDK]} token to make uncertainty a first-class output, and Ji et al.~\citep{ji2025calibratingverbaluncertaintylinear} identify a single linear ``verbal uncertainty'' direction that can be intervened on at inference time to reduce hallucinations. Wu et al.~\citep{wu2026mitigatingllmhallucinationbehaviorally} show that standard RLVR incentivises ``good test-takers'' over ``honest communicators,'' and propose proper-scoring-rule training that encourages stochastic abstention.

AbstentionBench~\citep{kirichenko2025abstentionbench} is the most directly related abstention benchmark: 20 datasets and roughly 35k unanswerable questions across 6 scenarios (unknown answers, underspecification, false premises, subjective interpretations, stale data, and reasoning-induced abstention). Critically, reasoning fine-tuning degrades abstention by 24\%, revealing a tension between capability and restraint. \textsc{AgentAbstain} extends this from text-only QA to tool-using agents, adding two dimensions: \emph{timing} (before vs.\ during execution) and \emph{side effects} (reversibility of actions).

\textbf{Over-refusal and calibration in the opposite direction.}
Over-abstention is equally problematic. OR-Bench~\citep{cui2025orbench} found high correlation between safety and over-refusal across 80k prompts, and R\"{o}ttger et al.~\citep{rottger2024xstest} (XSTest) earlier showed how superficially unsafe wording triggers refusal of benign queries. FalseReject~\citep{zhang2025falsereject} provides 16k benign-but-suspicious prompts across 44 categories with paired training signal. ORFuzz~\citep{zhang2025orfuzzfuzzingotherside} introduces evolutionary fuzzing for over-refusal discovery and an associated 1{,}786-case benchmark, while Yuan et al.~\citep{yuan2025overrefusalscenariobaseddiagnosticsposthoc} propose a scenario-based diagnostic and post-hoc mitigation framework. ACTOR~\citep{dabas2025justshiftsmitigatingoverrefusal} reduces over-refusal by 33\% via activation-level intervention. SORRY-Bench~\citep{xie2025sorrybench} found inconsistent safety across languages, and Campbell et al.~\citep{campbell2026defensiverefusal} showed LLMs refuse legitimate cybersecurity tasks at 43.8\%. RALMs are similarly affected: Zhou et al.~\citep{zhou2025retrievalaugmentedlanguagemodels} show that retrieval-augmented models exhibit characteristic over-refusal patterns tied to context confidence. These findings motivate our paired design, which measures both over- and under-abstention simultaneously.

\subsection{Clarification \& Ambiguity Handling}

A complementary line of work studies when LLMs should \emph{ask} rather than \emph{refuse}. CLAMBER~\citep{zhang2024clamberbenchmarkidentifyingclarifying} introduces a taxonomy of query ambiguity (contextual vs.\ LLM-oriented) and shows that current models struggle to ask high-quality clarifying questions even when they detect ambiguity. ClarQ-LLM~\citep{gan2024clarqllmbenchmarkmodelsclarifying} provides bilingual task-oriented dialog scenarios for evaluating clarification competence beyond yes/no questions. ClarifyMT-Bench~\citep{luo2025clarifymtbenchbenchmarkingimprovingmultiturn} reveals a consistent ``under-clarification bias'' across ten LLMs that worsens with dialogue depth, and proposes \textsc{ClarifyAgent} as a structured remedy. Noisy-ToolBench~\citep{wang2026learningaskllmagents} systematically perturbs tool-use queries with missing information, ambiguous references, factual inaccuracies, and tool-infeasibility, and finds that LLM agents tend to fabricate missing arguments rather than ask. AskBench~\citep{zhao2026askaskbenchrubricguidedrlvr} introduces rubric-guided RLVR specifically for the ask/answer decision, and HumanEvalComm~\citep{Wu_2025} measures clarification behaviour in code-generation agents. Our taxonomy's S1 (missing parameter) and S2 (ambiguous specification) scenarios align with this literature, but we additionally couple them with paired should-act variants and a runtime execution environment, allowing us to measure not just whether agents ask but whether they ask \emph{when it matters} for the action they would take.

\subsection{Agent Capability \& Tool-Use Benchmarks}

\textbf{General agent capability.}
AgentBench~\citep{liu2025agentbench} tests LLMs as agents across 8 environments. WebArena~\citep{zhou2024webarena} and BrowserGym~\citep{dechezelles2025browsergymecosystemwebagent} provide self-hosted web environments where state-of-the-art agents lag far behind humans (e.g., 78.24\% human vs.\ 14.41\% GPT-4 on WebArena). $\tau$-bench~\citep{yao2024taubench} introduced pass$^k$ for measuring consistency across multiple trials, and was followed by $\tau^2$-bench~\citep{barres2025tau2bench}, which adds a dual-control environment in which both user and agent can call tools. Terminal-Bench~\citep{merrill2026terminalbenchbenchmarkingagentshard} finds frontier models at $<$65\% on CLI tasks. FieldWorkArena~\citep{moteki2025fieldworkarena} and TheAgentCompany~\citep{xu2025theagentcompanybench} evaluate safety-critical and professional tasks respectively. 

\textbf{Tool-use specifically.}
BFCL V4~\citep{patil2025bfcl} is the most established tool-use leaderboard, evaluating serial/parallel function calls via AST matching and adding multi-turn and agentic categories; importantly, it includes a relevance-detection split that tests whether agents can recognise when no tool fits, measuring one slice of tool-selection abstention. ToolBeHonest~\citep{zhang2024toolbehonest} (ToolBH) formalises solvability detection, solution planning, and missing-tool analysis, and reports that frontier models score only 37--45/100 on solvability detection of unsolvable tool-use tasks. Critically, ToolBH evaluates agents \emph{statically}, asking the model to judge solvability over a tool set up front, whereas our setup requires \emph{dynamic execution}: agents must discover infeasibility at runtime through tool interaction. Furthermore, ToolBH covers only tool-level unsolvability (missing, potential, or limited tools), while we additionally test query-level (S1--S4) and runtime (S6--S8) abstention triggers. ToolSandbox~\citep{lu2025toolsandboxstatefulconversationalinteractive} introduces stateful tool execution and an ``insufficient information'' category that is closely related to our S1; we differ in coupling each insufficient-information case with a complete should-act counterpart. ToolMATH~\citep{choi2026toolmath} shows reasoning failures cascade through multi-tool pipelines, and OpaqueToolsBench~\citep{hallinan2026opaquetoolsbench} tests underspecified tools. ToolLLM~\citep{qin2023toolllm} provides 16,000 APIs for general-purpose tool-use training and evaluation. Most tool-use benchmarks score whether the agent reaches the goal; few systematically reward the agent for stopping when it should not.

\subsection{Agent Safety, Deception, \& Risk Benchmarks}

\textbf{Safety and harm refusal.}
AgentHarm~\citep{andriushchenko2025agentharm} finds 48--85\% refusal rates across 110 malicious agent tasks but only tests \emph{initial} refusal, not mid-task halting. Agent-SafetyBench~\citep{zhang2025agentsafetybench} covers 349 environments and 8 risk categories; no agent scores above 60\%. R-Judge~\citep{yuan2024rjudgebenchmarkingsafetyrisk} benchmarks safety risk awareness over agent interaction records across 27 scenarios. ToolEmu~\citep{ruan2024identifyingriskslmagents} formalises agent-tool interaction as a POMDP and uses LM emulation to surface high-stakes risks across 36 toolkits and 144 test cases without manual sandbox setup; our task model adopts a similar POMDP framing but couples it with a paired-task design rather than using LM-emulated grading. ST-WebAgentBench~\citep{levy2026stwebagentbench} introduces policy-aware metrics for web agents. SafeToolBench~\citep{xia2025safetoolbenchpioneeringprospectivebenchmark} evaluates tool-utilisation safety in a prospective rather than retrospective manner. 

\textbf{Deception and reward hacking.}
Guo et al.~\citep{guo2025agentsupwarddeceivers} introduce ``agentic upward deception'', the phenomenon in which an agent facing environmental constraints (broken tools, mismatched information sources) conceals failure and proceeds with unsupported actions. Their benchmark of 200 tasks across 8 realistic constrained scenarios reveals that 11 popular LLMs routinely guess results, perform unsupported simulations, substitute unavailable sources, and fabricate local files, and that prompt-based mitigation is only marginally effective. The trigger conditions they study (broken tools, mismatched sources) overlap substantially with our S6 (Critical Tool Failure) and S7 (Conflicting Evidence). The key difference is methodological: Guo et al.~\citep{guo2025agentsupwarddeceivers} detect deception \emph{post-hoc} by analysing trajectories on tasks that are \emph{designed} to elicit deception, whereas we evaluate abstention \emph{by design} via paired should-act and should-abstain tasks, allowing us to attribute behaviour to the abstention trigger $\delta$ rather than to baseline incompetence and to measure both over- and under-abstention in a single framework. Closely related, ImpossibleBench~\citep{zhong2025impossiblebenchmeasuringllmspropensity} mutates unit tests to create code tasks where any pass implies specification-violating shortcut behaviour; frontier models exhibit alarmingly high cheating rates (up to 76\% on Impossible-SWEbench). DeceptionBench~\citep{huang2025deceptionbenchcomprehensivebenchmarkai} provides 150 scenarios across five domains for measuring deceptive behaviour more broadly.

\subsection{The Gap}

Two gaps emerge from the landscape above. First, abstention research remains confined to single-turn textual QA; no prior benchmark extends it to tool-using agents whose halt decisions are sequential, grounded in environment state, and carry real side effects. Capability benchmarks measure completion, tool-use benchmarks rarely reward stopping, and clarification benchmarks measure asking but not the act/abstain decision under realistic execution. Second, agent benchmarks each isolate one slice of restraint (pre-execution clarification, static unsolvability, safety-driven refusal of malicious goals, or post-hoc deception detection) without jointly addressing \emph{when} a halt is required and \emph{what consequences} blind execution would produce. None evaluates abstention from legitimate tasks that become infeasible, epistemically unjustified, or harmful mid-execution: tasks that are not malicious, not adversarial, and would be entirely appropriate to act on under a slight variation of the inputs. \benchmark{} bridges both gaps as the first comprehensive benchmark for agentic abstention, unifying these fragmented evaluation axes along two dimensions: \emph{phase} (pre-execution vs.\ runtime triggers) and \emph{side effects} (informational vs.\ irreversible operational deliverables).


